\scanpage{59}
% @ lines 1–8 / text
Ici le signe $(*)$ en marge indique que l'équivalence des 3 catégories en question avec les 8 autres n'est valable que si $S$ est premier à $2$ -- en général on a seulemt des foncteurs qui vont de chacune des huit autres catégories dans chacune des trois marquées du signe $(*)$, mais ce ne sont pas des foncteurs pl. fidèles ....

% @ lines 9–16 / text
Un dictionnaire complet devrait comprendre, notamment, la description explicite des $9\times8$\inserted{$=72$} équivalences de catégories, sans compter les \deleted{$\unsure{11.10.9=990}$}\inserted{$9.8.7=504$} isomorphismes de transitivité (si on voulait être regardant ...) On fera un peu moins, mais on essaiera de donner une image claire de la situation.

% @ lines 17–21 / text
\uline{Principe général}. On va définir des représentants typiques des onze catégories dans le cas où $S_0=\operatorname{Spec}\mathbb Z$, et on définira une opération de $G_{0\deleted{\unsure{S}}}=\mathrm{GP}(1)$ sur chacun. On établira que ceci établit une isom.
% @ line 22 / display
\[
 \mathrm{GP}(1)\simeq\operatorname{Aut}(\unsure{\mathscr X|\mathcal X|\mathfrak X}_0)
\]
% @ lines 23–24, continues / text
de $\mathrm{GP}(1)$ avec le schéma des autom. de chacun des objets $\mathscr X$ envisagés -- \unsure{en se|on se|car si}
% archive sheet number: 58
